% [VERIFY] methodology -- solution/{algorithm,architecture,constraints,heuristics}.md, concepts.md, src/configs.
% Body = scientific contribution; code-level identifiers (residual_block.py, yaml keys) -> Appendix.
% TikZ residual-block schematic owed by venue-shaping rule (method's defining visual).
\section{Deep Residual Learning}
\label{sec:method}

\subsection{Residual Learning}
\label{sec:residual}
Let us consider $\Hmap(\xin)$ as an underlying mapping to be fit by a few
stacked layers, with $\xin$ denoting the inputs to the first of these
layers. If one hypothesizes that multiple nonlinear layers can asymptotically
approximate complicated functions, then it is equivalent to hypothesize that
they can asymptotically approximate the residual function, i.e.,
$\Hmap(\xin) - \xin$ (assuming that the input and output are of the same
dimensions). So rather than expecting stacked layers to approximate
$\Hmap(\xin)$, we explicitly let these layers approximate a residual function
$\Fmap(\xin) := \Hmap(\xin) - \xin$. The original function thus becomes
$\Fmap(\xin) + \xin$. Although both forms should be able to asymptotically
approximate the desired functions, as hypothesized, the ease of learning might
be different. This reformulation is motivated by the counterintuitive
\degrad{}: if the added layers can be constructed as identity mappings, a deeper
model should have training error no greater than its shallower counterpart. The
\degrad{} suggests that solvers might have difficulties in approximating
identity mappings by multiple nonlinear layers. With the residual learning
reformulation, if identity mappings are optimal, the solvers may simply drive
the weights of the multiple nonlinear layers toward zero to approach identity
mappings. We note that the hypothesis that stacked nonlinear layers can
asymptotically approximate the residual function remains an open
question~\cite{montufar2014}.

\subsection{Identity Mapping by Shortcuts}
\label{sec:shortcut}
We adopt residual learning to every few stacked layers. A building block is
defined as
\begin{equation}
\yout = \Fmap(\xin, \{W_i\}) + \xin .
\label{eq:block}
\end{equation}
Here $\xin$ and $\yout$ are the input and output vectors of the layers
considered. The function $\Fmap(\xin, \{W_i\})$ represents the residual mapping
to be learned. For the two-layer basic block illustrated in
\cref{fig:resblock}, $\Fmap = W_2\,\sigma(W_1\xin)$, in which $\sigma$ denotes
ReLU and the biases are omitted for simplifying notations. The operation
$\Fmap + \xin$ is performed by a shortcut connection and element-wise addition,
and we adopt the second nonlinearity after the addition, i.e.,
$\sigma(\yout)$. The shortcut connections in \cref{eq:block} introduce neither
extra parameter nor computation complexity, which is important for our
comparisons between \plainnet{}s and residual nets that have the same number of
parameters, depth, width, and computational cost.

The dimensions of $\xin$ and $\Fmap$ must be equal in \cref{eq:block}. When this
is not the case --- for example, when changing the input/output channels at a
stage boundary --- we perform a linear projection $W_s$ by the shortcut
connection to match the dimensions:
\begin{equation}
\yout = \Fmap(\xin, \{W_i\}) + W_s\,\xin .
\label{eq:proj}
\end{equation}
The form of the residual function $\Fmap$ is flexible, but a function with a
single layer is observed to provide no advantage, behaving similarly to a linear
layer; we therefore use a residual function $\Fmap$ with two or more layers.

\begin{figure}[t]
\centering
\resizebox{0.92\columnwidth}{!}{%
\begin{tikzpicture}[
  node distance=6mm,
  box/.style={draw, rounded corners, minimum width=30mm, minimum height=6mm,
              align=center, font=\footnotesize},
  op/.style={draw, circle, inner sep=1pt, font=\footnotesize},
  >={Stealth[length=2mm]}]
  \node (in) {\footnotesize $\xin$};
  \node[box, below=of in] (c1) {$3\times3$ conv};
  \node[box, below=of c1] (b1) {BN};
  \node[box, below=of b1] (r1) {ReLU};
  \node[box, below=of r1] (c2) {$3\times3$ conv};
  \node[box, below=of c2] (b2) {BN};
  \node[op, below=of b2] (add) {$+$};
  \node[box, below=of add] (r2) {ReLU};
  \node[below=5mm of r2] (out) {\footnotesize $\yout=\Fmap(\xin)+\xin$};
  \draw[->] (in) -- (c1);
  \draw[->] (c1) -- (b1);
  \draw[->] (b1) -- (r1);
  \draw[->] (r1) -- (c2);
  \draw[->] (c2) -- (b2);
  \draw[->] (b2) -- (add);
  \draw[->] (add) -- (r2);
  \draw[->] (r2) -- (out);
  % identity shortcut
  \draw[->, rounded corners] (in.east) -- ++(18mm,0)
        node[midway, above, font=\scriptsize]{identity $\xin$} |- (add.east);
\end{tikzpicture}}
\caption{Residual learning: a basic two-layer building block. The shortcut
connection performs identity mapping and its output is added element-wise to the
output of the stacked layers (\cref{eq:block}). Identity shortcuts add neither
parameters nor computational complexity. \emph{Source:
\texttt{architecture.md} (basic block); \texttt{algorithm.md} (Eqn.~1).}}
\label{fig:resblock}
\end{figure}

\subsection{Network Architectures}
\label{sec:arch}
We have tested various \plainnet{}s and residual nets, and have observed
consistent phenomena. The \plainnet{} baselines are mainly inspired by the
philosophy of VGG nets. The convolutional layers mostly have $3\times3$ filters
and follow two simple design rules: (i) for the same output feature map size,
the layers have the same number of filters; and (ii) if the feature map size is
halved, the number of filters is doubled so as to preserve the time complexity
per layer. We perform downsampling directly by convolutional layers that have a
stride of $2$. The network ends with a global average pooling layer and a
1000-way fully-connected layer with softmax. It is worth noticing that our model
has fewer filters and lower complexity than VGG nets: our 34-layer \plainnet{}
has $3.6$ billion FLOPs, which is only $18\%$ of VGG-19 ($19.6$ billion FLOPs).

Based on the \plainnet{}, we insert shortcut connections which turn the network
into its counterpart residual version. The identity shortcuts of
\cref{eq:block} can be used directly when the input and output are of the same
dimensions. When the dimensions increase, we consider two options: (A) the
shortcut still performs identity mapping, with extra zero entries padded for
increasing dimensions, introducing no extra parameter; or (B) the projection
shortcut of \cref{eq:proj} is used to match dimensions, done by $1\times1$
convolutions. For both options, when the shortcuts go across feature maps of two
sizes, they are performed with a stride of $2$. \Cref{tab:arch} details the
five-stage skeleton shared by all depths, with the per-stage block counts and
overall complexity.

\begin{table*}[t]
\centering
\caption{Architectures for ImageNet. Building blocks are shown in brackets, with
the numbers of blocks stacked. Downsampling is performed by conv3\_1, conv4\_1,
and conv5\_1 with a stride of 2. \emph{Source: Table 1
(\texttt{evidence/tables/table1\_imagenet\_architectures.md}); exact FLOPs and
block counts.}}
\label{tab:arch}
\small
\setlength{\tabcolsep}{4pt}
\begin{tabular}{l l l l l l}
\toprule
layer name & output size & 18-layer & 34-layer & 50-layer & 101/152-layer \\
\midrule
conv1   & $112\times112$ & \multicolumn{4}{l}{$7\times7$, 64, stride 2} \\
\midrule
\multirow{2}{*}{conv2\_x} & \multirow{2}{*}{$56\times56$}
  & \multicolumn{4}{l}{$3\times3$ max pool, stride 2} \\
\cmidrule(l){3-6}
  & & $\left[\begin{smallmatrix}3\times3,64\\3\times3,64\end{smallmatrix}\right]\!\times2$
    & $\left[\begin{smallmatrix}3\times3,64\\3\times3,64\end{smallmatrix}\right]\!\times3$
    & $\left[\begin{smallmatrix}1\times1,64\\3\times3,64\\1\times1,256\end{smallmatrix}\right]\!\times3$
    & $\left[\begin{smallmatrix}1\times1,64\\3\times3,64\\1\times1,256\end{smallmatrix}\right]\!\times3$ \\
\midrule
conv3\_x & $28\times28$
  & $\left[\begin{smallmatrix}3\times3,128\\3\times3,128\end{smallmatrix}\right]\!\times2$
  & $\left[\begin{smallmatrix}3\times3,128\\3\times3,128\end{smallmatrix}\right]\!\times4$
  & $\left[\begin{smallmatrix}1\times1,128\\3\times3,128\\1\times1,512\end{smallmatrix}\right]\!\times4$
  & $\left[\begin{smallmatrix}1\times1,128\\3\times3,128\\1\times1,512\end{smallmatrix}\right]\!\times\{4,8\}$ \\
\midrule
conv4\_x & $14\times14$
  & $\left[\begin{smallmatrix}3\times3,256\\3\times3,256\end{smallmatrix}\right]\!\times2$
  & $\left[\begin{smallmatrix}3\times3,256\\3\times3,256\end{smallmatrix}\right]\!\times6$
  & $\left[\begin{smallmatrix}1\times1,256\\3\times3,256\\1\times1,1024\end{smallmatrix}\right]\!\times6$
  & $\left[\begin{smallmatrix}1\times1,256\\3\times3,256\\1\times1,1024\end{smallmatrix}\right]\!\times\{23,36\}$ \\
\midrule
conv5\_x & $7\times7$
  & $\left[\begin{smallmatrix}3\times3,512\\3\times3,512\end{smallmatrix}\right]\!\times2$
  & $\left[\begin{smallmatrix}3\times3,512\\3\times3,512\end{smallmatrix}\right]\!\times3$
  & $\left[\begin{smallmatrix}1\times1,512\\3\times3,512\\1\times1,2048\end{smallmatrix}\right]\!\times3$
  & $\left[\begin{smallmatrix}1\times1,512\\3\times3,512\\1\times1,2048\end{smallmatrix}\right]\!\times3$ \\
\midrule
        & $1\times1$ & \multicolumn{4}{l}{average pool, 1000-d fc, softmax} \\
\midrule
\multicolumn{2}{l}{FLOPs}
  & $1.8\times10^{9}$ & $3.6\times10^{9}$ & $3.8\times10^{9}$
  & $7.6\times10^{9}\,/\,11.3\times10^{9}$ \\
\bottomrule
\end{tabular}
\end{table*}

\subsection{Deeper Bottleneck Architectures}
\label{sec:bottleneck}
For deeper nets, we modify the building block as a \emph{bottleneck} design out
of concern for the training time that we can afford. For each residual function
$\Fmap$, we use a stack of three layers instead of two. The three layers are
$1\times1$, $3\times3$, and $1\times1$ convolutions, where the $1\times1$ layers
are responsible for reducing and then restoring dimensions, leaving the
$3\times3$ layer a bottleneck with smaller input/output dimensions. This design
has a time complexity similar to that of the two-layer block. The
parameter-free identity shortcuts are particularly important for the bottleneck
architectures: if the identity shortcut is replaced with a projection, one can
show that the time complexity and model size are doubled, as the shortcut is
connected to the two high-dimensional ends. So identity shortcuts lead to more
efficient models for the bottleneck designs. We construct 50-, 101-, and
152-layer residual nets by using this bottleneck block, with the per-stage
block counts in \cref{tab:arch}. Notably, the 152-layer residual net has $11.3$
billion FLOPs, which is still of lower complexity than VGG-16/19 nets ($15.3$/$19.6$
billion FLOPs).

\subsection{Implementation}
\label{sec:impl}
For each residual block we follow the forward computation summarized in
\cref{alg:resblock}. The shortcut branch is selected once at construction time
according to whether the block changes dimensions and which shortcut option is
in effect; identity shortcuts dominate, with projections used only at the rare
stage boundaries. Our training recipe follows the established practice for deep
networks on ImageNet: batch normalization is applied right after each
convolution and before activation, weights are initialized with the MSRA scheme,
and the networks are trained from scratch by SGD. Exact optimizer settings,
augmentation, and test-time protocols are deferred to the Appendix so that the
body remains at the level of the architectural contribution.

\begin{algorithm}[t]
\caption{Forward pass of one residual block.}
\label{alg:resblock}
\begin{algorithmic}[1]
\Require feature $\xin$; residual function $\Fmap$ (basic or bottleneck);
         shortcut option; downsample flag
\State $\text{out} \gets \Fmap(\xin)$
  \Comment{conv$\to$BN$\to$ReLU$\to$conv$\to$BN (or $1{\times}1,3{\times}3,1{\times}1$)}
\If{$\dim(\text{out})=\dim(\xin)$ \textbf{and} not downsample}
   \State $\text{id} \gets \xin$ \Comment{parameter-free identity}
\Else
   \If{option $=$ A}
      \State $\text{id} \gets \textsc{ZeroPadChannels}(\textsc{Stride2}(\xin))$
   \Else
      \State $\text{id} \gets \textsc{Projection}(\xin)$
        \Comment{$1{\times}1$ conv (+stride 2)}
   \EndIf
\EndIf
\State $\text{out} \gets \text{out} + \text{id}$ \Comment{element-wise add}
\State \Return $\textsc{ReLU}(\text{out})$
\end{algorithmic}
\end{algorithm}

